\subsection{Express: Funktionalität}

\begin{frame}{Express: Funktionalitäten}{Applikationsprotokolle}
  Das Expressmodul ist in API und Funktion dem express Framework von node.js nachempfunden.
  \begin{itemize}
    \item \texttt{express}-"`Klasse"' von Node.JS ist \texttt{WebApp}-Klasse in SNode.C.
    \item \texttt{Router}-"`Klasse"' von Node.JS ist \texttt{Router}-Klasse in SNode.C
    \item Routing-API und -semantik mit \texttt{.get(\ldots)}, \texttt{.post(\ldots)}, \ldots, \texttt{.use(\ldots)}, \texttt{.all(\ldots)}, \texttt{next()}, \texttt{next("route")} und \texttt{next("router")} gleicht wie in Node.JS/Express
    \item Application- bzw. Middleware-Callbacks sind \alert{Lambda}-Expressions
    \item Behandlung von \alert{cookies}, \alert{queries} und \alert{header} gleich wie in node.js-express
    \item \alert{Regular-Expression} Route-Matching
    \item Aktuell existierende Middlewares:
    \begin{description}[BasicAuthentication Middleware]
      \item [Static Middleware] Auslieferung von statischen HTML-Sites
      \item [JSON Middleware] Parst JSON-Body in JSON-Datenstrucktur (nlohmann JSON C++ Library)
      \item [VHost Middleware] Unterstützt virtual Hosts
      \item [BasicAuthentication Middleware] Legacy Username/Password Authentifikation
    \end{description}
  \end{itemize}
\end{frame}
